\section{承接而不抹平：六个问题怎样彼此牵动}
\label{sec:synthesis-connections-costs}

迁徙、家户、制度、边境、宗教与文本不是六个可以各自写完的抽屉。一次迁徙会让户籍失准，也会改变寺院、市场和学校能够接触到谁；一项均田或兵役安排需要道路与地方中介，进而影响家庭能否维持劳作；一座边城的得失改变军粮和人群流动，又可能让宗教文本、工匠与城市记忆沿新的路线移动。综合写作的目标不是把这些关系压成一个“大势”，而是让读者在具体转折中看到哪一根线被拉紧，成本最终落在谁身上。

\subsection{第一根线：迁徙与行政互相制造}

人们常把迁徙看成战争的后果，登记看成战争之后的修复；两者实际上彼此制造。朝廷因失去人口、税源和兵员而要重新登记，登记的分类、差役与土地安排又会促使某些家户迁走、寄居或寻找新的保护。东晋侨置保存北方旧籍，使南渡者能够在新朝廷中主张名分；其后土断把实际住处重新纳入地方责任，又会改变同一家庭与旧籍之间的关系。这里没有一条从混乱到秩序的直线，只有在不同资源条件下反复调整的可见性。

北魏的均田与三长同样如此。国家试图借土地、家户和基层编组使人群可被追踪，迁徙、战争和隐匿则不断提醒簿册并不等于实况。制度能把一些原先不稳定的关系写入常规责任，也可能把未能登记、无地或依附关系复杂的人推向更不稳定的位置。将迁徙和行政放在一起，便不会把流民写成纯粹被动的数字，也不会把户籍写成天然中性的技术。

\subsection{第二根线：家庭劳动支撑军事选择}

北伐、守城、迁都和边防都需要粮食、布帛、舟车、马匹与维修。它们最后要由家庭劳动、地方市场和中介组织提供。江淮沿线的一次军事集结，可能要求船户提供运输、乡里交出粮食、工匠修造器具；关中或代北的一次动员，则受耕作周期、草场、山口与仓储限制。将军在地图上推进的距离，常常是后方家庭能承受多少役使的另一种表达。

这种联系也使军事失败的后果不止是撤军。粮食被消耗、劳力被调走、道路被破坏之后，家庭需要更长时间恢复；地方强者若能在此时提供借贷、保护或粮食，依附关系便可能加深。国家想把征发常规化，是为了避免每次危机都以掠夺方式取得资源；可常规赋役本身也会把风险固定在特定人群身上。制度得失应当同时问军队能否持续和家庭能否持续，不能只以某年胜负作答。

\subsection{第三根线：边境与城市相互定义}

边境看似远离都城，实际不断重塑城市。河西的绿洲、代北的军镇、淮汉的渡口和江陵的上游通道，把马匹、粮食、僧侣、文书、使者和逃亡者送往长安、平城、洛阳、邺或建康。城市越依赖某条道路，就越会把边地的安全、价格与人群移动带入自己的仓储和市场。反过来，城中的官署、寺院、作坊和军营又把对纸墨、木材、粮食与劳力的需求传回道路。

北魏迁都洛阳不是边境从此退场，而是中心重新计算边境成本。新都更接近中原腹地和礼制资源，旧军镇的给养、地位和政治联系却不能因此消失；六镇动乱显示这种距离会在特定条件下转化为危机。南朝以长江和淮河组织防卫时也一样：都城的安全依赖沿线船只、城镇和家户的配合。将边境写成“外部威胁”会掩盖它如何每天参与都城生活。

\subsection{第四根线：宗教网络既承受又重塑资源}

寺院与道观位于这张关系网的交点。它们需要供养、土地、建筑和人员，因而会受战争、迁徙和政策变化影响；它们又可以保存经卷、组织仪式、提供行旅节点和积累社会声望，使流动的人与物有机会在城市中停留。长安译场、洛阳寺院、河西石窟和建康的宗教网络之所以值得并读，就在于它们展示同一信仰并不以同一种资源形式存在。

禁佛政策最能检验这种复杂性。国家把僧尼、器物或寺产纳入清查，确实改变了许多人的位置，也可能为军政或财政取得资源；但寺院不是可以一次性清空的库房，信仰与文本也不是被诏令触及就完全消失的对象。地方执行、人员流动、私人供养和后来的恢复都使后果不均。把宗教仅写成统治者宽容或严酷的性格，会把这一整套资源与关系的重组误作一句道德判断。

\subsection{第五根线：文本把断裂保存成可争论的历史}

文本并不站在这些变化之外。诏令要把行政目标传下去，户籍把家户化为可追问的对象，墓志把迁徙家庭的身份钉在新旧地名之间，造像题记把供养者与地点联结，经卷与家训则让教育、信仰和处世跨过一次次政权转换。不同文本保存不同的尺度，正使我们能够从帝纪的事件线走向家庭、城市和边地。

文本也会把断裂组织成意义。正史常以一朝兴亡安排前后，文学会把失都、流亡和归隐编成可感的记忆，目录与地理书则把变化的世界重新分类。它们并不需要互相说同一种话。读者若知道一篇赋不是统计、一方墓志不是全社会、一项诏令不是地方普查，就能让这些文本共同扩大问题，而不让任何一种材料霸占结论。

\subsection{第六根线：再统一的收益与成本}

隋灭陈使南北对峙的最高政治框架结束，收益首先是行政、交通和征发可以在更大空间内被重新协调。沿江、淮和关中的道路不再必须服务于彼此敌对的朝廷，州郡接收也有机会减少反复的边界战争。但收益不是没有成本的：江南的家户、寺院、官员和市场须进入新的登记与调度，北方军政网络则要承担治理更广大区域的责任。统一能否持久，取决于这些成本是否被承认、分配和处理。

这也是两晋南北朝留给后世的实际问题。任何宣称天下归一的政权，都要回答怎样使流动的人被安置、怎样使家庭能在赋役下维持、怎样使边地和城市共享道路与供给、怎样处理宗教与知识网络、怎样让记录不只是中央的愿望。答案不可能只有一种，也不可能由一份诏书完成。把六根线并置，读者便能理解：王朝更替不是宏大结构压过个人生活，而是在无数具体关系中不断试探其限度。

\begin{causalchain}[综合章的读法：从一件事回到六个问题]
遇到一场战争，可问它使哪些道路、粮食和家庭劳动紧张；遇到一项改革，可问它怎样重新登记迁徙者并依赖地方中介；遇到寺院营造或禁令，可问资源、城市和文本网络怎样变化；遇到一座城的易手，可问边境与市场怎样重新连结。六个问题不会给出同一答案，却能防止我们把任何一次胜败写成没有社会后果的名称。
\end{causalchain}

因此，本卷的综合并非最后一页的总结词。它是一种回读方法：从长安的译场回到关中战争与书写资源，从洛阳迁都回到军镇与家庭的距离，从建康的北伐回到水运和侨籍，从河西的城镇回到绿洲与文本，从隋的受降回到接收一座城市和一批家户的细节。时间线由此没有被主题取代，反而获得了更长的因果回声。
